% Nouveau EGA V (cote 25): archive cote 25, site page 26 (PDF page 27).
% typescript. Draft compared: none.
% AI-transcribed by Claude Opus 5.5 (claude-opus-5-5, adaptive thinking, effort high), Message Batches API, batch ega25b (msgbatch_01A7n5YuXubUPNK1zsepi8RG), 2026-11-13; not yet checked by eye.
\scanpage{25/26}
\mspage{22}
enfin que X ne peut être disconnecté par une partie fermée Z de X telle
que dim $\overline{f(Z)} \unsure{\leq} 0$.    Sous ces conditions, pour tout $\hand{\zeta \in} \overset{\hand{\vee}}{P}$ , $Y_{\hand{\zeta}}$ est
géométriquement connexe.

Démonstration.  Comme \xout{xesxf(X} aucun des $f(X_i)$ n'est fini, on voit
que chacune d es composantes irréductibles $\underline{Y}_i$ de $\underline{Y}$ domine $\overset{\hand{\vee}}{P}$ ;
d'autre part $\underline{Y} \to \overset{\hand{\vee}}{P}$ est propre (car $\underline{Y}$ est propre sur k, l'étant sur X
qui l'est sur k).  D'autre part, en vertu de (3.3.) la fibre générique
$Y_{\hand{\eta}}$ de $\underline{Y} \to \overset{\hand{\vee}}{P}$ est géométr\xout{i}iquement connexe. Enfin, $\overset{\hand{\vee}}{P}$ est régulier et
à fortiori géom unibranche. Il suffit maintenant d'appliquer 15.6.3.
(\xout{xsxxxxxxxxxxxxxxxxxxxxxxxx}\inserted[typed]{qui n'est qu'une variante du} th. de connexité  de Zariski) pour concl
re que \uline{toutes} les fibres  de $\underline{Y} \to \overset{\hand{\vee}}{P}$ sont géom connexes, cqfd.

\margin{$|$} En fait,  il n'est pas difficile, par une démsnstration de type an
logue, de généraliser \over{\illegible{1}}{5}.1. dans le même sens que dans 4.4. Si tu ne
veux pas te taper l'exercice toi-même, signale au, moins la chose en re
marque. Y dire aussi que l'on ne fait aucune discrimination dans \over{\illegible{1}}{5}.1.
à l'égard des sections hyperplanes qui auraient une dimension excéden-%
\margin{$|$} taire.  Du point de vue plan, il sera peut-être plus clair de bloquer
toutes les questions de connexité (incluant 3.3. et 4.4. ) dans un
même n$^{\text{o}}$ .
\deleted{6. \uline{Application à la construction de sections hyperplanes} \hand{/}\inserted[typed]{\uline{ou de multisections}}\hand{)} \uline{de type préci-%
sé}, \xout{ouxdexmultisectionsx}}
\mspage{26}
